\subsection{Simultaneous caps on the ten primitive sources}
\begin{statement}{P3F.45}{Proved}
Use any of the ten actual primitive-translation sources of P3K.22,
with its marked line, translation $G$, and selected simple whole
$I_1$ zero. Install finitely many actual root-line caps away from
that zero. Retain all elliptic, line and annular coefficients.
After the exact finite-kernel comparison of P3K.19, write their
effective norm pairs as $(m_i,n_i)$, with $m_i>0$ and
$\gcd(m_i,n_i)=1$.
At a multiplicative site require the actual component-transitive
root; at an additive site use the complete local domains of
P3F.42--43. Thus the additive possibilities for these primitive
sources are reduced III of order $4l+1$, reduced IV of order
$3l+1$, and type II with reduced order $6l+1$ or good order
$6l$, the latter with $l\geq1$. Every source has at most one
such additive site.

For the circle orbit $Y$, the complete degree-two gluing map has
cokernel
\[
 \operatorname{coker}D_2
   =\bigoplus_{i\text{ smooth or multiplicative}}
                         (\mathbb Z/m_i)^2.             \tag{F45.1}
\]
Its kernel is zero. This cokernel injects into $H_2(Y;\mathbb Z)$.
Consequently every smooth or multiplicative effective order of a
sphere member is one. Among the actual constructions in this
statement the sphere criterion is exactly this order restriction
together with
\[
 \begin{cases}
 n_*+m_*\displaystyle\sum_{i\ne *}n_i=\pm1,
                         &\text{one additive site }*,\\[2mm]
 \displaystyle\sum_i n_i=\pm1,
                         &\text{no additive site}.
 \end{cases}                                           \tag{F45.2}
\]
In that sector the full fundamental group is cyclic of the
displayed order, with $\mathbb Z$ for order zero. The criterion
holds for every permitted one-vertex cusp fan. It retains the
actual order-one source and line clutches; no equality of their
parameter returns is asserted by this topology calculation.
\end{statement}
\begin{proof}
\par\medskip\noindent\textbf{The ordinary comparison and integral exterior bases.}\enspace
The relative cusp product of P3F.9 identifies the orbit before
these replacements with the actual mapping torus
$W_0=T_{t_{-G}}S$, with $(p,1)\sim(p-G,0)$ and the physical
outward gamma marking. The integral calculation in P3K.22 gives
$H_2(W_0)=\mathbb Z[O]$ and $H_3(W_0)=\mathbb ZF$.
Let $q\geq1$ be the number of replaced tubes, $W$ their exterior,
$U_i$ the ordinary tubes, and $E_i=\partial U_i$.
All their central source components have multiplicity one, and
$G$ acts transitively on them. The ordinary tube's degree-three
component trace maps to the primitive class $F$ of $W_0$.
At a smooth fibre this trace is the ordinary three-torus.
The sum of these maps is therefore onto $H_3(W_0)$.
The actual section $O$ has boundary $\sum_i\alpha_i$ in
$\bigoplus H_1(E_i)$, where each $\alpha_i$ is its primitive
positive base meridian. This is a nonzero primitive vector.
The full ordinary Mayer--Vietoris sequence now gives the integral
isomorphism
\[
 \bigoplus_iH_2(E_i)
       \ \xrightarrow{\ (i_2,-j_2^{\rm ord})\ }\
 H_2(W)\oplus\bigoplus_iH_2(U_i).                    \tag{F45.3}
\]
In particular $H_2(W)$ is free, and the images of the ordinary
cap kernels form an integral basis of it.

For clarity, all the other exterior groups used below have the
following actual source-pair description. Let $L_i$ be the columns
of the source component classes of the $i$th fibre, and let $B$
be the integral source Gram matrix already computed in P3K.22.
Writing $L=(L_1,\ldots,L_q)$, the relative source complex has
\[
 C_4=\mathbb Z,\quad C_3=\mathbb Z^q,\quad
 C_2=\operatorname{NS}(S),\quad C_1=\mathbb Z^{\sum r_i},
 \quad C_0=\mathbb Z,
 \qquad d_4=\mathbf1,\quad d_2=L^tB.                 \tag{F45.4}
\]
The other differentials are zero. The boundary inclusions are
$L_i$ in degree two, their component-coordinate embeddings in
degree one, and the $i$th coordinate in degree three. These are
the actual pair maps: the source and the tube retractions have
free even homology, so their split even complexes have no odd
derived homotopy parameter. This argument includes the tangencies
of III and IV. The primitive trivial lattice of the verified
torsion-free Mordell--Weil source makes the cokernel of $d_2$
free of rank $q-1$. Its generators are the base meridians with
their single sum relation. Translation acts identically on them
and on $H_3$ of the source exterior. The mapping-torus cone of
(F45.4) consequently gives
\[
 H_1(W)=\mathbb Z^q,\qquad H_4(W)=\mathbb Z^{q-1},
 \qquad H_3(W)\text{ is free}.                       \tag{F45.5}
\]
The stated primitivity can also be read directly from the
integral bases in P3K.22: adjoining $O$ and $G$ to the fibre and
the nonidentity components gives the full source lattice, with
the actual Mordell--Weil quotient $\mathbb Z$.

\par\medskip\noindent\textbf{The actual inversion and every degree-two column.}\enspace
On $W_0$ the map $I[p,t]=[-p,-t]$ is an involution. It preserves
the base and every chosen tube, reverses the five-dimensional
orientation, and preserves the normal base-disc orientation.
On a smooth tube choose the actual Lie lift $G(q)$ and the map
\[
 (a,b,t)\longmapsto[a+b\tau(q)+tG(q),t].
\]
In these coordinates $I$ is exactly $(-a,-b,-t)$ and fixes the
base meridian $\alpha$. Thus $ab,a\gamma,b\gamma$ are even,
whereas $a\alpha,b\alpha,\alpha\gamma$ are odd.
The same formula on a Tate sector descends through
$a\mapsto a+Nb+c_Gt$. In the exact basis $(z,R,B,L)$ of
P3F.40, $z,B$ are even and $R,L$ are odd. The ordinary cap kernel
is precisely $\mathbb ZR\oplus\mathbb ZL$.

At a reduced additive ordinary site, $H_2(E_i)$ is free of
rank two. Write $F=r z$, with $r=1,2,3$ for II, III, IV.
The primitive class $z$ maps to the primitive component-orbit
class $C$ in the ordinary cap. The boundary is a three-torus
mapping torus and is parallelizable: a path from its oriented
linear monodromy to the identity trivializes its vertical
tangent bundle. Its rank-two unimodular intersection form is
therefore even. Choose $w$ with $z\cdot w=1$ and put
$s=w-(w^2/2)z$. The boundary kernel is isotropic, hence equals
$\mathbb Zs$. Inversion fixes $z$, preserves this kernel, and
reverses the intersection form. The equality
$Iz\cdot Is=-z\cdot s$ gives $Is=-s$. No $z$ correction is
possible in the preserved cap kernel.

Every ordinary kernel vector in (F45.3) is thus odd. Hence
$I=-1$ on the entire free group $H_2(W)$. Every even boundary
core column has zero exterior image, because twice that image
is zero. This proves their simultaneous integral vanishing.
The map $I$ is used on the original exterior throughout.

\par\medskip\noindent\textbf{The full local replacement maps.}\enspace
At a smooth site retain
\[
 \mu=m\alpha+n\gamma+\lambda_a a+\lambda_b b,
 \qquad mu+nv=1.
\]
In cap basis $(a,b,h)$ the actual first-homology quotient is
\[
 Q=\begin{pmatrix}
 1&0&-u\lambda_a&-v\lambda_a\\
 0&1&-u\lambda_b&-v\lambda_b\\
 0&0&-n&m
 \end{pmatrix}                                      \tag{F45.6}
\]
on $(a,b,\alpha,\gamma)$. Its kernel is exactly
$\mathbb Z\mu$. The higher maps are its exterior powers,
because the actual smooth root and extending translation give
the marked product $T^3\times D^2$ as in P3K.22.
In degree-two order
$(ab,a\gamma,b\gamma,a\alpha,b\alpha,\alpha\gamma)$,
the complete local block is
\[
 D_2^{\rm sm}=
 \begin{pmatrix}
 0&0&0&1&0&0\\0&0&0&0&1&0\\0&0&0&0&0&1\\
 -1&v\lambda_b&-v\lambda_a&u\lambda_b&-u\lambda_a&0\\
 0&-m&0&n&0&\lambda_a\\
 0&0&-m&0&n&\lambda_b
 \end{pmatrix}.                                    \tag{F45.7}
\]
Eliminating the first three exterior rows and the unit core
column leaves the two diagonal entries $m,m$.
In degree-three order
$(N=ab\alpha,F\gamma=ab\gamma,a\alpha\gamma,b\alpha\gamma)$,
the cap row is
\[
                    (-n,m,-\lambda_b,\lambda_a).     \tag{F45.8}
\]

At a multiplicative site use exactly the marked data of P3F.40,
including $c=-c_G$, $A=Nj-cn$, and $\gcd(m,A)=1$.
Choose $p,r$ with $mp+nr=-j$. In degree one, on
$(\alpha,b,\gamma)$, the actual cap map is
\[
       \begin{pmatrix}p&1&r\\-n&0&m\end{pmatrix}.
\]
The complete degree-two block on $(z,R,B,L)$ is
\[
 D_2^{\rm mult}=
 \begin{pmatrix}
 0&1&0&0\\0&0&0&1\\-m&A&0&0\\0&0&-m&-A
 \end{pmatrix}.                                    \tag{F45.9}
\]
This retains the actual integral cap basis change proved by
boundary isotropy in P3F.40. Its cokernel is $(\mathbb Z/m)^2$.
For degree three use $(N=ab\alpha,F\gamma=ab\gamma,
V=a\gamma\alpha)$. The intersection pairing with
$(\alpha,b,\gamma)$ is
\[
 \begin{pmatrix}0&-1&0\\0&0&-1\\1&0&0\end{pmatrix}.
\]
The primitive meridian bounds in the cap. Poincare--Lefschetz
duality identifies the cap degree-three row with its negative
dual, with the sign fixed by the actual fibre trace
$F\gamma\mapsto m$. Thus this row is exactly
\[
                             (-n,m,j).                \tag{F45.10}
\]
Primitivity follows also from the cap's component-trace
generator; there is no further index in this duality map.

For the actual additive roots of P3F.42--43, the cap has
$H_1=H_2=H_3=\mathbb Z$. Reduced III/IV specialize the fibre
to $rC$ and have boundary $F=r z$, hence $z\mapsto C$.
Reduced II has the same assertion with $r=1$.
For good II the actual smooth-root transport has degree one,
and $D-1$ is unimodular on the elliptic lattice; again
$z\mapsto C$. With its remaining actual collar coefficient
$\kappa$, the block on $(z,s)$ is
\[
 D_2^{\rm add}=\begin{pmatrix}0&1\\-1&-\kappa\end{pmatrix},
 \qquad j_1=j_3=(-n,m)                             \tag{F45.11}
\]
in the corresponding $(\alpha,\gamma)$ and $(N,F\gamma)$
bases. These specialization statements use the actual resolved
root and line maps, not the numerical Kodaira symbols alone.
The simultaneous integral exterior basis makes (F45.7),
(F45.9), and (F45.11) independent diagonal blocks of $D_2$.
They are injective and prove (F45.1).

\par\medskip\noindent\textbf{Every remaining gluing degree.}\enspace
Put $c_i=3,2,1$ at smooth, multiplicative, additive sites.
In $H_1(W)$ use the $q-1$ independent base meridians and the
common gamma generator. Elliptic boundary loops map to zero;
$\alpha_q=-\sum_{i<q}\alpha_i$, and all $\gamma_i$ have the
same image. These statements follow from (F45.4) and from the
actual source-slice paths, so they retain the original markings.
Together with (F45.6), (F45.10)'s degree-one precursor, and
(F45.11), they specify every entry of $D_1$.

The following integral description likewise specifies $D_3$.
For an ordinary cap its kernel in degree three is generated by
$N_i$ and $c_i-1$ local elliptic-rim classes. All these classes
are inversion-even, by the coordinates just used. The source
fundamental class of $W_0$ has boundary $\sum_iN_i$.
The ordinary degree-three exact sequence, whose cap traces
surject onto the primitive $F$, shows that the images of these
kernel classes have this one primitive relation and no other
relation. Adjoin one image of $F\gamma$ to obtain an integral
basis of $H_3(W)$. Every other $F\gamma$ image equals it:
their differences lie in the even kernel and are also odd.
The freeness in (F45.5) makes those differences zero.
Consequently the exterior columns are $N_q=-\sum_{i<q}N_i$,
one common $F\gamma$, and the independent local rim columns.
The complete cap rows are (F45.8), (F45.10), and (F45.11).

Eliminating the local elliptic columns in $D_1$ and the local
rim rows in $D_3$ gives identity blocks and the same $2q$ by
$2q$ matrix $S_q$. Its columns are
$(\alpha_1,\gamma_1,\ldots,\alpha_q,\gamma_q)$, or
$(N_1,F\gamma_1,\ldots,N_q,F\gamma_q)$, respectively.
Its first $q-1$ rows have $1$ in column $\alpha_i$ and $-1$
in column $\alpha_q$. Its next row has $1$ in every gamma
column. Its last $q$ rows have $(n_i,-m_i)$ in the $i$th
pair and zero elsewhere. For example,
\[
 S_2=\begin{pmatrix}
 1&0&-1&0\\0&1&0&1\\n_1&-m_1&0&0\\0&0&n_2&-m_2
 \end{pmatrix}.                                    \tag{F45.12}
\]
Expanding after eliminating the $q-1$ meridian rows gives
\[
 \det S_q=(-1)^{q(q+1)/2}
       \sum_i n_i\prod_{j\ne i}m_j.                 \tag{F45.13}
\]
Equivalently, over the rational function field eliminate each
cap row to set $\gamma_i=(n_i/m_i)\alpha_i$; the remaining
meridian relation has coefficient $\sum_i n_i/m_i$.
Multiplication by $\prod_i m_i$ and the displayed column
permutation gives (F45.13) as a polynomial identity over
$\mathbb Z$, with no restriction on the signed $n_i$.
Finally
\[
 D_4=(I_{q-1}\mid-\mathbf1),\qquad
 D_0=\begin{pmatrix}\mathbf1^t\\-I_q\end{pmatrix}. \tag{F45.14}
\]
Thus every nonzero degree of the full Mayer--Vietoris map is
specified as an integer matrix, including every local collar.
The program \path{code/part3_root/primitive_multicap_complete.py}
exports them, checks the exterior-power and duality identities,
and executes 121 full integral cases with positive, negative and
zero determinant. The symbolic determinants and the integral
eliminations prove the assertions for all orders.

When $D=\sum_i n_i\prod_{j\ne i}m_j\ne0$, the complete
sequence in particular gives
\[
 \begin{gathered}
 H_0(Y)=H_5(Y)=\mathbb Z,\qquad H_4(Y)=0,\\
 H_1(Y)=H_3(Y)=\operatorname{coker}S_q,\qquad
 H_2(Y)=\bigoplus_{i\text{ smooth or multiplicative}}
                                  (\mathbb Z/m_i)^2.
 \end{gathered}                                    \tag{F45.15}
\]
For a sphere all the latter orders are therefore one. Eliminating
their cap generators in $S_q$ leaves $D=n_*+m_*\sum_{i\ne *}n_i$,
or $D=\sum_i n_i$ with no additive site. Its Smith factors are
units followed by $|D|$; for $D=0$ the final factor is zero.

\par\medskip\noindent\textbf{The full fundamental group and the sphere implication.}\enspace
On the complement of the source critical values the section
gives the actual group extension by the elliptic lattice.
Its normal elliptic subgroup is abelian. Every multiplicative
replacement kills its original primitive vanishing circle:
transport to the actual root central polygon collapses precisely
that circle. This statement includes its graph translation.
A reduced additive replacement kills the elliptic subgroup by
its simply connected reduced central fibre. At a good II cap,
the boundary elliptic lattice maps onto the full elliptic
normal subgroup of the cap. The other source fibres already
kill that lattice. Indeed the complement of the sole II fibre
has base a disc; its fundamental group is the quotient of the
elliptic lattice by the retained vanishing cycles. Its integral
first homology is zero from (F45.4) with $L=F$, so that abelian
group is zero. With no additive replacement, the original
vanishing-cycle relations of the simply connected rational
source are all retained in the unchanged fibres or in the
multiplicative caps. This again kills the entire lattice.
These are relations in the actual group, and every source
translation acts trivially on that elliptic normal subgroup.

After eliminating it, van Kampen therefore gives precisely
\[
 \left\langle\gamma,\alpha_i,h_i\ \left|\
 \prod_i\alpha_i=1,\quad [\alpha_i,\gamma]=1,\quad
 h_i^{m_i}=\gamma,\quad h_i^{-n_i}=\alpha_i
                         \right.\right\rangle .       \tag{F45.16}
\]
All elliptic path coefficients have disappeared here because
their full subgroup has just been killed, not because a
different boundary marking was substituted. With every smooth
and multiplicative $m_i=1$, each corresponding $h_i$ equals
$\gamma$. There is at most one other generator $h_*$, and
$\gamma=h_*^{m_*}$. The sole remaining relation is
$h_*^{n_*+m_*\sum_{i\ne *}n_i}=1$. This proves the cyclic
claim and the exact signs in (F45.2). Without an additive site
the relation is $\gamma^{\sum_i n_i}=1$.

For $D=\pm1$, (F45.7)--(F45.14) give the integral homology
of $S^5$ and (F45.16) gives simple connectivity. The resulting
oriented orbit is the standard $S^5$. The actual delta circle
is free in every added cap, and its single fixed sphere at the
selected cusp has the primitive positive Hopf meridian. The
Gysin and fixed-tube maps of P3K.22 then give simple connectivity
and the integral homology of $S^6$ for $X$. The vanishing of
$\Theta_6$ gives its standard oriented diffeomorphism type.
Conversely the orbit calculation of P3F.7 applied to a sphere
member forces these same orbit groups to vanish, giving
(F45.1)--(F45.2). The actual common blowup and finite flop chain
of P3F.35 transports this topology criterion to every allowed
cusp fan. All local finite characters, base powers and
order-one gluing maps remain those of the constructed family.
\end{proof}
